% =====================================================================
\section{Model 2: Learning from Data and Optimal Control}
\label{sec:model2}
% =====================================================================

Model~1 needs the physics: the mass $m$, the inertia $J$, the rotor constants. Model~2 asks a different question:
\emph{if we only had recorded flight data, could we still build a good controller?} It learns a linear model from data
(blocks 1--2) and then computes the best controller for that model in two ways (\cref{fig:flow2}):
\begin{itemize}
  \item \textbf{LQR} (blocks 3--6): the mathematically best \emph{linear} controller. Very cheap to run, but it knows
        nothing about limits.
  \item \textbf{MPC} (blocks 7--9): plans the next second optimally \emph{and} respects limits such as
        ``never faster than 1\,m/s'', at the cost of solving a small optimisation problem every step.
\end{itemize}
Block 10 flies the real, nonlinear drone with the result. All Model~2 matrices are in \emph{discrete time} with
$\Delta t=0.04$\,s (25\,Hz, the \code{control\_rate}): $\vx_{k+1}=A\vx_k+B\vu_k$.
Toy model files:
\begin{itemize}
  \item \code{M2\_1\_Flight\_Data\_DMDc.mlx} (blocks 1--2),
  \item \code{M2\_2\_LQR\_Bellman\_Riccati\_DARE.mlx} (blocks 3--6),
  \item \code{M2\_3\_MPC\_QP\_Apply\_to\_Drone.mlx} (blocks 7--10).
\end{itemize}

% ---------------------------------------------------------------------
\subsection{Block 1 --- Recording flight data}
\label{sec:m2b1}
% ---------------------------------------------------------------------

\begin{idea}
Fly the drone and write down, 25 times per second, where it was and what we told it to do. To make the data useful we
measure everything as a \emph{difference from hover}, and we deliberately give the drone small random pushes so that the
data contains enough variety to learn from.
\end{idea}

\begin{eqbox}{Equation (M2.1): sampled state and input, as deviations from hover}
\begin{equation}
  \vx_k=\begin{bmatrix}\vp(t_k)-\vp_{\mathrm{ref}}\\ \vv(t_k)\\ \Log\big(R(t_k)\big)\\ \vw(t_k)\end{bmatrix}\in\R^{12},\qquad
  \vu_k=\begin{bmatrix}T(t_k)-mg_0\\ \vtau(t_k)\end{bmatrix}\in\R^{4},\qquad t_k=k\,\Delta t . \tag{M2.1}
\end{equation}
\end{eqbox}

\subsubsection*{Why it is written this way}
\textbf{Differences from hover.} Hover is a balance point of (1)--(4): level, still, thrust $=$ weight, every rate is
zero. Near it, the linear model $\vx_{k+1}=A\vx_k+B\vu_k$ has \emph{no constant term}. If we used the raw thrust $T$
instead of $T-mg_0$, gravity would appear as a constant $+mg_0$ that this form cannot represent, and the fit would be poor.

\textbf{Attitude as three numbers.} The learning method needs plain vectors, so the rotation matrix is replaced by its
rotation vector $\Log(R)$ (\cref{sec:maths}). Near hover these three numbers are roll, pitch and yaw in radians.

\textbf{A pilot plus random pushes.} A free drone is unstable (\cref{sec:m1b3}), so the data must be recorded with some
stabilising pilot flying it (the role ArduPilot plays). But if the input were \emph{only} feedback, $\vu=-K_0\vx$, every
input would be a fixed combination of the states, and the computer could not tell ``what the drone does'' apart from
``what the pilot does''. Adding small random pushes ($\pm2$\,N of thrust, $\pm0.03$\,N\,m of torque, each held for 0.2\,s) fixes
this. This is called \textbf{persistent excitation}.

\begin{toy}[the recorded data set]
Five flights of 25\,s each at 25\,Hz, starting from random points near $[0,0,2]$\,m: $5\times625$ samples of
$\vx_k\in\R^{12}$ and $\vu_k\in\R^4$, with a little sensor noise. Flights 1--4 are used to learn, flight 5 is kept back to
test. The largest tilt in all the data is 6.5$^\circ$, well inside the range where the drone is nearly linear (\cref{sec:m1b2}).
\end{toy}

\begin{figure}[H]
  \centering
  \includegraphics[width=0.55\linewidth]{M2_1_Flight_Data_DMDc_1.png}
  \caption{The five recorded flights around the hover point.}
\end{figure}

\begin{toycode}{M2\_1\_Flight\_Data\_DMDc}{Block 1}
Xr(:,k) = DroneLib.to_dev(x, P.p_hover) + noise;        % logged state x_k
if mod(k-1, 5) == 0                                       % new random push
    exc = [2; 0.03; 0.03; 0.01].*(2*rand(4,1) - 1);       %   every 0.2 s
end
[~, ~, u_app] = DroneLib.mixer(u_trim + DroneLib.pilot(Xr(:,k), P) + exc, P);
Ur(:,k) = u_app - u_trim;                                % logged input u_k
\end{toycode}

\begin{real}
In the project, block 1 is a ROS~2 bag recording of \code{/ap/v1/pose/filtered}, \code{/ap/v1/twist/filtered} and the
commands we send, from SITL or a real flight (\code{ros2 bag record ...}). The pose and twist are exactly the inputs our
node already reads, so no new sensor is needed, and nothing about mass or inertia needs to be known.
\end{real}

% ---------------------------------------------------------------------
\subsection{Block 2 --- DMDc: learning $A$ and $B$ from the data}
\label{sec:m2b2}
% ---------------------------------------------------------------------

\begin{idea}
Look at thousands of moments of the form ``the drone was \emph{here}, we pushed like \emph{this}, and 40\,ms later it was
\emph{there}'', and find the single linear rule $\vx_{k+1}=A\vx_k+B\vu_k$ that explains all of them best. It is curve
fitting for motion. It is the EDMD idea from the course notes (\code{K = Psi2*pinv(Psi1)}) with the input added; the
name DMDc means \emph{Dynamic Mode Decomposition with control}.
\end{idea}

\begin{eqbox}{Equations (M2.2)--(M2.4): snapshot matrices and the least-squares model}
\begin{align}
  X=\big[\vx_1\ \cdots\ \vx_{N-1}\big],\quad X^+&=\big[\vx_2\ \cdots\ \vx_{N}\big],\quad U=\big[\vu_1\ \cdots\ \vu_{N-1}\big] \tag{M2.2}\\
  \big[\,A\ \ B\,\big]&=X^+\begin{bmatrix}X\\ U\end{bmatrix}^{\dagger} \tag{M2.3}\\
  \vx_{k+1}&\approx A\,\vx_k+B\,\vu_k \tag{M2.4}
\end{align}
\end{eqbox}

\subsubsection*{Derivation}
Put the states side by side in a matrix $X$, the same states shifted one step later in $X^+$, and the inputs in $U$. We want
one matrix $G=[A\ \ B]$ such that, for every column,
\[
\vx_{k+1}=A\vx_k+B\vu_k=G\begin{bmatrix}\vx_k\\ \vu_k\end{bmatrix}.
\]
Written for all columns at once:
\[
X^+\approx G\,\Omega,\qquad \Omega=\begin{bmatrix}X\\U\end{bmatrix}.
\]
With noisy data this cannot hold exactly, so we choose the $G$ with the smallest total squared prediction error
$\|X^+-G\Omega\|^2$ (least squares, \cref{sec:lsq}). Setting the derivative with respect to $G$ to zero:
\[
-2\,(X^+-G\Omega)\,\Omega^\top=0\quad\Longrightarrow\quad G\,\Omega\Omega^\top=X^+\Omega^\top
\quad\Longrightarrow\quad G=X^+\Omega^\top(\Omega\Omega^\top)^{-1}=X^+\Omega^\dagger .
\]
$\Omega\Omega^\top$ is a $16\times16$ matrix. It can only be inverted if $\Omega$ has 16 independent rows ($12+4$), which is
exactly what the random pushes of block 1 guarantee. The first 12 columns of $G$ are $A$ and the last 4 are $B$.

\textbf{What should the answer be?} If the data came from the physics of Model~1, the exact answer is the discrete
version of the Jacobians \eqref{eq:AB}: $A=e^{A_c\Delta t}$ and $B=\int_0^{\Delta t}e^{A_cs}ds\,B_c$ (the ``zero-order hold'':
the input is held constant for one step). So we can check DMDc against the physics.

\begin{toy}[DMDc by hand, one number]
The true system is $x_{k+1}=0.9x_k+0.5u_k$, but we pretend not to know $0.9$ and $0.5$. Three recorded steps:
$X=[1,\ 1.4,\ 0.76]$, $U=[1,\ -1,\ 0.5]$, $X^+=[1.4,\ 0.76,\ 0.934]$. Then
\[
\Omega\Omega^\top=\begin{bmatrix}1+1.96+0.5776 & 1-1.4+0.38\\ 1-1.4+0.38 & 1+1+0.25\end{bmatrix}=\begin{bmatrix}3.5376&-0.02\\-0.02&2.25\end{bmatrix},
\]
\[
X^+\Omega^\top=\begin{bmatrix}1.4+1.064+0.70984 & 1.4-0.76+0.467\end{bmatrix}=\begin{bmatrix}3.17384&1.107\end{bmatrix},
\]
\[
[a\ \ b]=\frac{1}{3.5376\cdot2.25-0.02^2}\begin{bmatrix}3.17384&1.107\end{bmatrix}\begin{bmatrix}2.25&0.02\\0.02&3.5376\end{bmatrix}=\begin{bmatrix}0.9000&0.5000\end{bmatrix}.
\]
The hidden numbers are found exactly, because this data has no noise.

\textbf{The drone} (2500 snapshots, rank of $\Omega$ = 16):
\begin{center}\small
\begin{tabular}{@{}llrr@{}}
\toprule
entry & physical meaning & learned by DMDc & physics\\
\midrule
$A_{1,4}$ & $\Delta t$ (position from velocity) & 0.0399 & 0.0400\\
$A_{4,8}$ & $g_0\Delta t$ ($v_x$ from pitch) & 0.3909 & 0.3924\\
$A_{5,7}$ & $-g_0\Delta t$ ($v_y$ from roll) & $-0.3862$ & $-0.3924$\\
$B_{6,1}$ & $\Delta t/m$ ($v_z$ from thrust) & 0.0268 & 0.0267\\
$B_{10,2}$ & $\Delta t/J_x$ (roll rate from roll torque) & 2.0007 & 2.0000\\
$B_{12,4}$ & $\Delta t/J_z$ (yaw rate from yaw torque) & 1.0096 & 1.0000\\
\bottomrule
\end{tabular}
\end{center}
Overall the learned $A$ is within 2.6\% of the physics and $B$ within 1.1\%. The data even ``knows'' the mass:
$\Delta t/B_{6,1}=1.494$\,kg (true: 1.5\,kg). On flight 5, which was never used for learning, predicting 1\,s ahead from the
inputs alone gives a position error of 1.33\,cm, against 1.26\,cm for the exact physics model.
\end{toy}

\begin{figure}[H]
  \centering
  \includegraphics[width=0.5\linewidth]{M2_1_Flight_Data_DMDc_2.png}
  \caption{Eigenvalues of the learned $A$ (crosses) and of the physics (circles). All lie at or very near 1: a
  chain of integrators, exactly as block 3 of Model 1 predicted.}
\end{figure}

\begin{toycode}{M2\_1\_Flight\_Data\_DMDc}{Block 2}
Omega = [Xa; Ua];
rank_Omega = rank(Omega)                % must be 12 + 4 = 16 (enough excitation)
G = Xpa*pinv(Omega);                    % (M2.3)
A = G(:, 1:12);
B = G(:, 13:16);
\end{toycode}

\begin{real}
DMDc runs offline, on a laptop, on a recorded bag. If the drone changes (a heavier camera, a new battery, a different
frame), a two-minute flight is enough to learn the new $A$ and $B$, without measuring anything by hand.
\end{real}

% ---------------------------------------------------------------------
\subsection{Block 3 --- The Bellman equation: what is ``the best'' controller?}
\label{sec:m2b3}
% ---------------------------------------------------------------------

\begin{idea}
Before we can find the best controller we must say what ``best'' means. We give the flight a \emph{score}: every moment we
are off target costs points, and every push of the motors costs points. The best controller is the one with the lowest
total score. Bellman's idea: the best plan from here is the best \emph{first move} plus the best plan from wherever that
move takes us.
\end{idea}

\begin{eqbox}{Equations (M2.5)--(M2.7): cost, Bellman equation, quadratic value function}
\begin{align}
  J&=\sum_{k=0}^{\infty}\big(\vx_k^\top Q\,\vx_k+\vu_k^\top R\,\vu_k\big) \tag{M2.5}\\
  V(\vx_k)&=\min_{\vu_k}\Big[\vx_k^\top Q\vx_k+\vu_k^\top R\vu_k+V(\vx_{k+1})\Big],\qquad \vx_{k+1}=A\vx_k+B\vu_k \tag{M2.6}\\
  V(\vx)&=\vx^\top P\,\vx \tag{M2.7}
\end{align}
\end{eqbox}

\subsubsection*{(M2.5) The score}
$\vx^\top Q\vx$ is a weighted sum of squared errors and $\vu^\top R\vu$ a weighted sum of squared pushes. Big $Q$: ``I hate
being off target''. Big $R$: ``I hate using the motors''. We choose them with \textbf{Bryson's rule}: $Q_{ii}=1/x_{i,\max}^2$
and $R_{jj}=1/u_{j,\max}^2$, where $x_{i,\max}$ is the largest error we are willing to accept:
\begin{center}\small
\begin{tabular}{@{}lll@{}}
\toprule
quantity & largest acceptable value & weight\\
\midrule
position error & 0.5\,m & $1/0.5^2=4$\\
velocity & 1\,m/s (\code{max\_velocity}) & 1\\
roll, pitch / yaw & 0.2\,rad / 0.5\,rad & 25 / 4\\
angular rates & 1\,rad/s & 1\\
extra thrust & 5\,N & 0.04\\
roll, pitch torque / yaw torque & 0.2\,N\,m / 0.05\,N\,m & 25 / 400\\
\bottomrule
\end{tabular}
\end{center}
So a 50\,cm position error ``hurts'' as much as flying at 1\,m/s, tilting 0.2\,rad, or using 5\,N of extra thrust.

\subsubsection*{(M2.6) Principle of optimality}
Split the infinite sum into the first step and all the rest:
\[
V(\vx_k)=\min_{\vu_k,\vu_{k+1},\dots}\Big[\ell(\vx_k,\vu_k)+\sum_{j>k}\ell(\vx_j,\vu_j)\Big]
=\min_{\vu_k}\Big[\ell(\vx_k,\vu_k)+\underbrace{\min_{\vu_{k+1},\dots}\sum_{j>k}\ell(\vx_j,\vu_j)}_{V(\vx_{k+1})}\Big],
\]
where $\ell$ is the score of one step. Whatever the first move is, the remaining moves must be the best ones for the place
it leads to; otherwise we could improve them and lower the total. A problem with infinitely many unknowns becomes a
problem with \emph{one} unknown, $\vu_k$.

\subsubsection*{(M2.7) The value function is a bowl}
Guess that the lowest possible score from $\vx$ is a quadratic, $V(\vx)=\vx^\top P\vx$. Then the bracket in (M2.6) is a
quadratic in $\vu_k$. Its minimum (block 4) is linear in $\vx_k$, and putting it back gives again a quadratic in $\vx_k$
(block 5). The guess is self-consistent, so all that is left to find is the matrix $P$.

\begin{toy}[one Bellman step by hand]
System $x_{k+1}=1.1x_k+0.5u_k$ (on its own it grows 10\% per step: unstable). $q=r=1$, the next value is $V(x)=x^2$, and
we are at $x=1$:
\[
\text{score}(u)=\underbrace{1}_{qx^2}+\underbrace{u^2}_{ru^2}+\underbrace{(1.1+0.5u)^2}_{V(x_{k+1})},\qquad
\frac{d}{du}=2u+(1.1+0.5u)=0\ \Rightarrow\ u^*=-\frac{1.1}{2.5}=-0.44,
\]
\[
V(1)=1+0.1936+(1.1-0.22)^2=1.968 .
\]
MATLAB's \code{fminsearch} on the same function returns $u^*=-0.4400$ and $1.9680$.
\end{toy}

\begin{toycode}{M2\_2\_LQR\_Bellman\_Riccati\_DARE}{Block 3}
a = 1.1;  b = 0.5;  q = 1;  r = 1;
cost = @(u) q*1^2 + r*u.^2 + 1*(a*1 + b*u).^2;
u_star = -a*b/(r + b^2)                 % by hand
[u_search, V1] = fminsearch(cost, 0)   % by the inbuilt optimiser
\end{toycode}

% ---------------------------------------------------------------------
\subsection{Block 4 --- The Linear Quadratic Regulator (LQR)}
\label{sec:m2b4}
% ---------------------------------------------------------------------

\begin{idea}
Do the minimisation of block 3 once, in general, with matrices. The answer is beautifully simple: the best push is a fixed
matrix $K$ times the current error. The whole controller is one $4\times12$ table of numbers.
\end{idea}

\begin{eqbox}{Equations (M2.8)--(M2.9): optimal gain and control law}
\begin{align}
  K&=\big(R+B^\top PB\big)^{-1}B^\top PA \tag{M2.8}\\
  \vu_k&=-K\,\vx_k \tag{M2.9}
\end{align}
\end{eqbox}

\subsubsection*{Derivation}
Put $V(\vx_{k+1})=(A\vx+B\vu)^\top P(A\vx+B\vu)$ into the bracket of (M2.6). It is a bowl in $\vu$ (\cref{sec:minquad}).
Its slope with respect to $\vu$, using $\partial(\vu^\top M\vu)/\partial\vu=2M\vu$, must be zero:
\[
2R\vu+2B^\top P(A\vx+B\vu)=\bm0\quad\Longrightarrow\quad(R+B^\top PB)\,\vu=-B^\top PA\,\vx .
\]
$R+B^\top PB$ is positive definite (because $R$ is), so it can be inverted and the point really is the bottom of the bowl.
Hence $\vu^*=-(R+B^\top PB)^{-1}B^\top PA\,\vx=-K\vx$.

\begin{toy}[the scalar system, and the drone's gain]
With the final $P=3.1215$ (from block 6): $K=\dfrac{ab\,P}{r+b^2P}=\dfrac{1.1\cdot0.5\cdot3.1215}{1+0.25\cdot3.1215}=0.9643$.
With control, $x_{k+1}=(a-bK)x_k=0.6179\,x_k$. Without control the state grew 10\% per step; with LQR it
\emph{shrinks} 38\% per step (\cref{sec:eig}).

\textbf{The drone's $K$}, two of its 48 numbers (from block 6):
\[
\delta T=-9.04\,\delta z-6.92\,v_z,\qquad
\tau_x=0.273\,\delta y+0.321\,v_y-1.444\,\theta_x-0.276\,\omega_x .
\]
Read them: too high or rising $\Rightarrow$ less thrust. Too far north ($\delta y>0$) $\Rightarrow$ positive roll torque, which
tips the drone towards $-y$, back to the target; but less so if it is already rolled ($\theta_x$) or rolling ($\omega_x$).
\end{toy}

% ---------------------------------------------------------------------
\subsection{Block 5 --- The Riccati recursion: finding $P$ step by step}
\label{sec:m2b5}
% ---------------------------------------------------------------------

\begin{idea}
$K$ needs $P$, and $P$ is ``the lowest score from here''. Start with a controller that looks 0 steps ahead, then 1 step,
then 2, and so on. Each time, the plan sees one more step of the future. After enough steps, looking further does not
change anything: we have the infinite-horizon $P$.
\end{idea}

\begin{eqbox}{Equation (M2.10): Riccati difference equation}
\begin{equation}
  P_{k+1}=A^\top P_kA-A^\top P_kB\big(R+B^\top P_kB\big)^{-1}B^\top P_kA+Q,\qquad P_0=Q \tag{M2.10}
\end{equation}
($k$ counts ``steps of look-ahead'', not time.)
\end{eqbox}

\subsubsection*{Derivation}
Put the best push $\vu^*=-K\vx$ back into the bracket of (M2.6):
\[
V_{\text{new}}(\vx)=\vx^\top\big[Q+K^\top RK+(A-BK)^\top P(A-BK)\big]\vx .
\]
Multiply out $(A-BK)^\top P(A-BK)=A^\top PA-A^\top PBK-K^\top B^\top PA+K^\top B^\top PBK$ and add $K^\top RK$. From (M2.8),
$(R+B^\top PB)K=B^\top PA$, so $K^\top(R+B^\top PB)K=K^\top B^\top PA$, and two terms cancel:
\[
V_{\text{new}}=\vx^\top\big[Q+A^\top PA-A^\top PBK\big]\vx=\vx^\top\big[Q+A^\top PA-A^\top PB(R+B^\top PB)^{-1}B^\top PA\big]\vx .
\]
The matrix in brackets is the new $P$: that is (M2.10).

\begin{toy}[the scalar recursion: $a=1.1$, $b=0.5$, $q=r=1$, $P_0=1$]
\[
P_{k+1}=1+1.21P_k-\frac{(0.55P_k)^2}{1+0.25P_k}
\]
\begin{center}\small
\begin{tabular}{l*{10}{c}}
\toprule
$k$ & 0 & 1 & 2 & 3 & 4 & 5 & 6 & 7 & 8 & $\infty$\\
$P_k$ & 1 & 1.968 & 2.596 & 2.905 & 3.036 & 3.089 & 3.109 & 3.117 & 3.120 & 3.1215\\
\bottomrule
\end{tabular}
\end{center}
$P_1=1.968$ is exactly the Bellman step of block 3. For the drone, starting from $P_0=Q$, the change becomes smaller than
$10^{-12}$ after 236 iterations (\cref{fig:riccati}): looking more than about 10\,s ahead changes nothing.
\end{toy}

\begin{figure}[H]
  \centering
  \includegraphics[width=0.55\linewidth]{M2_2_LQR_Bellman_Riccati_DARE_1.png}
  \caption{Riccati recursion for the drone: the change of $P$ falls steadily until it has converged.}
  \label{fig:riccati}
\end{figure}

\begin{toycode}{M2\_2\_LQR\_Bellman\_Riccati\_DARE}{Riccati recursion}
Pk = Q;
for it = 1:5000
    Kk = (R + B'*Pk*B) \ (B'*Pk*A);                       % (M2.8)
    Pn = A'*Pk*A - A'*Pk*B*Kk + Q;                        % (M2.10)
    change(it) = norm(Pn - Pk, 'fro')/norm(Pn, 'fro');
    Pk = Pn;
    if change(it) < 1e-12, break, end
end
\end{toycode}

% ---------------------------------------------------------------------
\subsection{Block 6 --- The DARE: jumping straight to the answer}
\label{sec:m2b6}
% ---------------------------------------------------------------------

\begin{idea}
When the recursion has stopped changing, $P_{k+1}=P_k$. Writing that down gives one equation for the final $P$, the
\emph{Discrete Algebraic Riccati Equation}. It can be solved directly instead of iterating. Its solution gives the final
gain $K$ and also tells us the total score in advance.
\end{idea}

\begin{eqbox}{Equation (M2.11): the Discrete Algebraic Riccati Equation}
\begin{equation}
  P=A^\top PA-A^\top PB\big(R+B^\top PB\big)^{-1}B^\top PA+Q \tag{M2.11}
\end{equation}
\end{eqbox}

\subsubsection*{Two useful facts}
\textbf{The total score is known in advance.} Rearranged, (M2.11) says $P=Q+K^\top RK+(A-BK)^\top P(A-BK)$. Along the
controlled flight $\vx_{k+1}=(A-BK)\vx_k$ this means
$\vx_k^\top P\vx_k-\vx_{k+1}^\top P\vx_{k+1}=\vx_k^\top Q\vx_k+\vu_k^\top R\vu_k$. Adding this up over all steps, everything in
between cancels (like $(a-b)+(b-c)+(c-d)=a-d$) and, because $\vx_k\to0$:
\[
J^*=\vx_0^\top P\,\vx_0 .
\]
\textbf{Solving without a toolbox.} The toy model solves (M2.11) with a standard trick (stable eigenvectors of a
$24\times24$ matrix pencil, \code{DroneLib.dare}), so no Control System Toolbox is needed.

\begin{toy}[the scalar DARE is a quadratic equation]
$P=1+1.21P-\dfrac{0.3025P^2}{1+0.25P}$. Multiply by $(1+0.25P)$ and collect terms: $0.25P^2-0.46P-1=0$, so
\[
P=\frac{0.46+\sqrt{0.46^2+1}}{0.5}=3.1215
\]
(the positive root), exactly the limit of the recursion.

\textbf{The drone.} The DARE solution satisfies (M2.11) to $6.5\times10^{-14}$ and agrees with the 236-step recursion to
$1.6\times10^{-11}$. The learned drone on its own has spectral radius $\rho(A)=1.017>1$ (unstable); with LQR
$\rho(A-BK)=0.950<1$ (stable). Starting 50/30/20\,cm away, the simulated total score is 41.4373, and the prediction
$\vx_0^\top P\vx_0$ is also 41.4373.
\end{toy}

\begin{figure}[H]
  \centering
  \includegraphics[width=0.55\linewidth]{M2_2_LQR_Bellman_Riccati_DARE_2.png}
  \caption{LQR on the learned model: the position errors die out smoothly.}
\end{figure}

\begin{toycode}{M2\_2\_LQR\_Bellman\_Riccati\_DARE}{DARE solved directly}
P = DroneLib.dare(A, B, Q, R);                     % (M2.11)
K = (R + B'*P*B) \ (B'*P*A);                       % (M2.8)
spectral_radius_closed = max(abs(eig(A - B*K)))   % < 1 : LQR makes it stable
\end{toycode}

\begin{real}[LQR]
On the drone, LQR would be one line every 40\,ms: $\vu=\vu_{\text{trim}}-K\vx$, with $\vx$ from the state estimate. That
costs 48 multiplications, nothing for any flight computer. Its weakness: it knows nothing about limits. Our node's
velocity law $\vv_{\text{cmd}}=0.8\ve_p+0.2\ve_v$ (\cref{sec:realctrl}) is a hand-tuned gain of exactly this kind, and it
needs a separate clip to respect \code{max\_velocity}.
\end{real}

% ---------------------------------------------------------------------
\subsection{Block 7 --- Model Predictive Control (MPC): planning with limits}
\label{sec:m2b7}
% ---------------------------------------------------------------------

\begin{idea}
MPC is like a chess player. Every 40\,ms it imagines the next second of flight, finds the best sequence of moves that never
breaks a rule (``never faster than 1\,m/s'', ``never ask a motor for more than it can give''), plays only the first move,
and then thinks again with the new position.
\end{idea}

\begin{eqbox}{Equation (M2.12): the finite-horizon problem with limits}
\begin{equation}
\begin{aligned}
  \min_{\vu_0,\dots,\vu_{N-1}}\ \ &\sum_{k=0}^{N-1}\big(\vx_k^\top Q\vx_k+\vu_k^\top R\vu_k\big)+\vx_N^\top P_f\,\vx_N\\
  \text{subject to}\ \ &\vx_{k+1}=A\vx_k+B\vu_k,\quad \vx_0=\text{the measured state},\\
  &\vu_{\min}\le\vu_k\le\vu_{\max},\qquad |C\vx_k|\le\vx_{\max}
\end{aligned}\tag{M2.12}
\end{equation}
\end{eqbox}

\subsubsection*{Reasoning}
Real drones have limits: total thrust between about 5 and 25\,N, the torque the rotors can produce, \code{max\_velocity}
= 1\,m/s, and a maximum tilt for a camera. With limits there is no neat formula like LQR's any more, so MPC solves the
problem numerically, and to keep it small:
\begin{itemize}
  \item it only plans $N$ steps ahead ($N=25$: one second);
  \item the infinite rest of the flight is replaced by a final cost $\vx_N^\top P_f\vx_N$.
\end{itemize}
\textbf{Why $P_f=P$ from the DARE.} We showed that $\vx_N^\top P\vx_N$ is \emph{exactly} the best score of the rest of the
flight when no limit is active. So with this choice, when no limit is touched, MPC gives exactly the LQR answer.
That is the link between the two branches of \cref{fig:flow2}.

The drone's limits in the toy model: $5\le T\le25$\,N (so $\delta T\in[-9.715,\,10.285]$\,N), $|\tau_{x,y}|\le0.3$\,N\,m,
$|\tau_z|\le0.05$\,N\,m, $|v_i|\le1$\,m/s, roll and pitch $\le0.35$\,rad (20$^\circ$). $C$ simply picks these five states out of $\vx$.

% ---------------------------------------------------------------------
\subsection{Block 8 --- Turning MPC into a Quadratic Program (QP)}
\label{sec:m2b8}
% ---------------------------------------------------------------------

\begin{idea}
Write every future state in terms of the unknown future pushes. Then the whole problem becomes: ``find the lowest point of
a bowl, but stay inside a fence''. That is called a \emph{quadratic program}, and computers solve it reliably in
milliseconds.
\end{idea}

\begin{eqbox}{Equations (M2.13)--(M2.15): stacked prediction, cost and constraints}
\begin{align}
  \mathbf X=S_x\vx_0+S_u\mathbf U,\qquad
  S_x=\begin{bmatrix}A\\ A^2\\ \vdots\\ A^N\end{bmatrix},\quad
  S_u=\begin{bmatrix}B&0&\cdots&0\\ AB&B&\cdots&0\\ \vdots&&\ddots&\vdots\\ A^{N-1}B&\cdots&AB&B\end{bmatrix} \tag{M2.13}\\[4pt]
  \min_{\mathbf U}\ \tfrac12\mathbf U^\top H\mathbf U+\bm f^\top\mathbf U,\qquad
  H=2\big(S_u^\top\bar QS_u+\bar R\big),\quad \bm f=2S_u^\top\bar QS_x\vx_0 \tag{M2.14}\\[2pt]
  \text{subject to}\quad G\mathbf U\le\bm h,\qquad A_{eq}\mathbf U=\bm b_{eq} \tag{M2.15}
\end{align}
$\mathbf X=[\vx_1;\dots;\vx_N]$, $\mathbf U=[\vu_0;\dots;\vu_{N-1}]$, $\bar Q=\diag(Q,\dots,Q,P_f)$, $\bar R=\diag(R,\dots,R)$.
\end{eqbox}

\subsubsection*{Derivation}
\textbf{(M2.13) Prediction.} Apply the model again and again:
$\vx_1=A\vx_0+B\vu_0$; $\vx_2=A\vx_1+B\vu_1=A^2\vx_0+AB\vu_0+B\vu_1$; and in general
$\vx_k=A^k\vx_0+\sum_{j=0}^{k-1}A^{k-1-j}B\vu_j$. Stacking $k=1..N$ gives (M2.13). The future is \emph{linear} in the unknown
pushes.

\textbf{(M2.14) Cost.} The score is $\vx_0^\top Q\vx_0+\mathbf X^\top\bar Q\mathbf X+\mathbf U^\top\bar R\mathbf U$. Substitute $\mathbf X$
and multiply out:
\[
\mathbf U^\top\big(S_u^\top\bar QS_u+\bar R\big)\mathbf U+2\,\vx_0^\top S_x^\top\bar QS_u\,\mathbf U+\underbrace{\vx_0^\top(Q+S_x^\top\bar QS_x)\vx_0}_{\text{does not depend on }\mathbf U}.
\]
Drop the last part (it does not change where the minimum is) and write the rest as $\tfrac12\mathbf U^\top H\mathbf U+\bm f^\top\mathbf U$.
$H$ is positive definite because $\bar R$ is, so the bowl has \emph{one} bottom and no false minima.

\textbf{(M2.15) Constraints.} Input limits: $\mathbf U\le\mathbf U_{\max}$ and $-\mathbf U\le-\mathbf U_{\min}$. State limits
$|\bar C\mathbf X|\le\bar{\vx}_{\max}$ become, after substituting $\mathbf X$, $\bar CS_u\mathbf U\le\bar{\vx}_{\max}-\bar CS_x\vx_0$ and
$-\bar CS_u\mathbf U\le\bar{\vx}_{\max}+\bar CS_x\vx_0$. Stacked, these give $G\mathbf U\le\bm h$. Equalities $A_{eq}\mathbf U=\bm b_{eq}$
can force, for example, ``be exactly at the target after $N$ steps''.

\textbf{The solver.} The toy model includes its own QP solver (\code{DroneLib.qp}, an interior-point method), so no
Optimization Toolbox is needed. For the drone it needs about 10 iterations.

\begin{toy}[a two-step QP by hand: $a=1.1$, $b=0.5$, $q=r=1$, $N=2$, $P_f=3.1215$, $x_0=1$]
\[
S_x=\begin{bmatrix}1.1\\1.21\end{bmatrix},\quad S_u=\begin{bmatrix}0.5&0\\0.55&0.5\end{bmatrix},\quad
\bar Q=\diag(1,\,3.1215),\quad \bar R=I,
\]
\[
H=\begin{bmatrix}4.3885&1.7168\\1.7168&3.5607\end{bmatrix},\qquad \bm f=\begin{bmatrix}5.2547\\3.7770\end{bmatrix}.
\]
\begin{itemize}
\item \textbf{No limits:} $\mathbf U^*=-H^{-1}\bm f=[-0.9643,\ -0.5958]$. The first move equals the LQR move $-Kx_0=-0.9643$
      \emph{exactly}, as promised by $P_f=P$.
\item \textbf{With $|u_k|\le0.5$:} $\mathbf U^*=[-0.5,\ -0.5]$. Both moves sit on the fence.
\item \textbf{Forcing $x_2=0$} with $[0.55\ \ 0.5]\,\mathbf U=-1.21$: $\mathbf U^*=[-1.3057,\ -0.9837]$. Check:
      $1.21+0.55(-1.3057)+0.5(-0.9837)=0$.
\end{itemize}
\textbf{The drone} ($N=25$): 100 unknowns and 450 inequalities. Starting 5\,cm off, no limit is touched and MPC gives
$[-0.4557,\,-0.0138,\,-0.0137,\,-0.0004]$, the same as $-K\vx_0$ to every printed digit. Starting 2\,m too low, LQR would ask for
$+18.1$\,N of extra thrust, far more than the allowed $+10.285$\,N. MPC uses exactly $+10.285$\,N and plans a climb that
touches, but never exceeds, 1\,m/s.
\end{toy}

\begin{toycode}{M2\_3\_MPC\_QP\_Apply\_to\_Drone}{Block 8, toy example}
Sx = [a; a^2];
Su = [b 0; a*b b];
Qbar = diag([1 Pf]);  Rbar = eye(2);
H = 2*(Su'*Qbar*Su + Rbar)                      % (M2.14)
f = 2*Su'*Qbar*Sx*x0
U_free    = (-H\f)'                             % no limits
U_limited = DroneLib.qp(H, f, [eye(2); -eye(2)], 0.5*ones(4,1))'   % |u| <= 0.5
\end{toycode}

% ---------------------------------------------------------------------
\subsection{Block 9 --- Receding horizon: plan, act a little, plan again}
\label{sec:m2b9}
% ---------------------------------------------------------------------

\begin{eqbox}{Equation (M2.16): receding-horizon control law}
\begin{equation}
  \mathbf U^*(\vx_k)=\big[\vu_0^*,\ \vu_1^*,\ \dots,\ \vu_{N-1}^*\big]
  \quad\Longrightarrow\quad \text{apply only }\vu_0^*,\ \text{measure }\vx_{k+1},\ \text{solve again} \tag{M2.16}
\end{equation}
\end{eqbox}

\begin{idea}
The plan is perfect only for the model, and the model is never perfect (it was learned, and the wind is not in it). So we
never follow a whole plan. We use its first move for 40\,ms, measure where we really are, and plan again. This turns a plan
into feedback: mistakes are corrected 25 times per second instead of piling up.
\end{idea}

\begin{toy}[8 s on the learned model, MPC against LQR]
Start 1.5\,m, 1.0\,m and 1.2\,m away from the target, at rest. MPC never exceeds 1\,m/s (largest speed 1.000\,m/s). LQR,
even when its pushes are clipped to the same motor limits, reaches 1.172\,m/s, because clipping the push does not limit
the \emph{speed}. The median time to solve one QP was 2.6\,ms, well inside the 40\,ms budget.
\end{toy}

\begin{figure}[H]
  \centering
  \includegraphics[width=0.6\linewidth]{M2_3_MPC_QP_Apply_to_Drone_1.png}
  \caption{Largest speed component over time: only MPC respects \code{max\_velocity} = 1\,m/s.}
\end{figure}

\begin{toycode}{M2\_3\_MPC\_QP\_Apply\_to\_Drone}{Block 9}
for k = 1:Kn
    um = DroneLib.mpc_solve(mpc, Xm(:,k));            % solve the QP, keep u_0*
    ul = min(max(-K*Xl(:,k), u_min), u_max);          % LQR, clipped
    Xm(:,k+1) = A*Xm(:,k) + B*um;
    Xl(:,k+1) = A*Xl(:,k) + B*ul;
end
\end{toycode}

% ---------------------------------------------------------------------
\needspace{14\baselineskip}
\subsection{Block 10 --- Applying it to the drone}
\label{sec:m2b10}
% ---------------------------------------------------------------------

\begin{eqbox}{Equations (M2.17)--(M2.18): from optimal input to motor speeds}
\begin{align}
  \vu_k&=\vu_{\text{trim}}+\begin{cases}\vu_0^*(\vx_k)&\text{(MPC)}\\ -K\vx_k&\text{(LQR)}\end{cases}
   =\begin{bmatrix}T_k\\ \vtau_k\end{bmatrix},\qquad \vu_{\text{trim}}=\begin{bmatrix}mg_0\\ \bm0\end{bmatrix} \tag{M2.17}\\
  \vf&=\Gamma^{-1}\vu_k,\qquad \Omega_i=\sqrt{f_i/k_f}\qquad(\text{the mixer of block 8 of Model 1}) \tag{M2.18}
\end{align}
\end{eqbox}

\subsubsection*{Reasoning}
The controllers of Model~2 work with \emph{differences from hover}. The motors need real values, so the hover thrust
$mg_0$ is added back. The rest is the same motor mixer as Model~1 (\cref{sec:m1b8}). The loop closes by measuring the new
state of the \emph{real, nonlinear} drone and running block 9 again.

\begin{toy}[one control sample, then a whole flight]
The drone is at $[1.5,\,-1.0,\,0.8]$\,m, at rest; the target is $[0,0,2]$. MPC gives
$\delta\vu=[10.285,\,-0.279,\,-0.300,\,0.002]$: full allowed climb thrust and the largest allowed roll and pitch torques. Adding
hover: $\vu=[25.000,\,-0.279,\,-0.300,\,0.002]$, rotor thrusts $\vf=[7.115,\,5.327,\,6.539,\,6.019]$\,N and motor speeds
$\vOmega=[912,\,789,\,875,\,839]$\,rad/s.

\textbf{The real test}: the nonlinear drone, with gusts and sensor noise, for 10\,s, flown by controllers designed only on
the model learned from data:
\begin{center}\small
\begin{tabular}{lrrr}
\toprule
controller & largest speed [m/s] & error after 10\,s [m] & motor speeds [rad/s]\\
\midrule
MPC & 0.994 & 0.015 & 584 -- 913\\
LQR (clipped) & 1.118 & 0.031 & 613 -- 911\\
\bottomrule
\end{tabular}
\end{center}
Both bring the real drone to the target. Only MPC keeps \code{max\_velocity} = 1\,m/s.
\end{toy}

\begin{figure}[H]
  \centering
  \includegraphics[width=0.6\linewidth]{M2_3_MPC_QP_Apply_to_Drone_2.png}
  \caption{The nonlinear drone flown by MPC and by LQR, both designed only from recorded data.}
\end{figure}

\begin{real}[Model 2]
Model~2 is not inside our ROS~2 node today; it is the data-driven alternative we studied and tested in the toy model. It
plugs into the same place as the SO(3) controller: it takes the state from \code{/ap/v1/pose/filtered} and
\code{/ap/v1/twist/filtered} (converted to $\vx_k$ with (M2.1)) and produces a command. Because ArduPilot accepts velocities,
the practical first step would be an MPC whose output is the velocity $\vv_{\text{cmd}}$ with the limit $\|\vv\|\le1$\,m/s written
into the problem, instead of the clip in \code{publish\_velocity\_command()}.
\end{real}

\subsection*{Recap of Model 2}
\begin{center}\small
\begin{tabularx}{\linewidth}{@{}c >{\raggedright\arraybackslash}p{2.6cm} >{\raggedright\arraybackslash}p{2.7cm} >{\raggedright\arraybackslash}p{2.5cm} X@{}}
\toprule
\textbf{\#} & \textbf{Block} & \textbf{In} & \textbf{Out} & \textbf{The idea in one sentence}\\
\midrule
1 & Flight data & drone, random pushes & $\{\vx_k\},\{\vu_k\}$ & Record what the drone did and what we told it, as differences from hover.\\
2 & DMDc & $X$, $X^+$, $U$ & $A$, $B$ & Least squares finds the linear rule that best predicts the next sample.\\
3 & Bellman & $A,B,Q,R$ & $V=\vx^\top P\vx$ & Best plan = best first move + best plan from where it leads.\\
4 & LQR & $A,B,R,P$ & $K$ & The bottom of a quadratic bowl is a linear law $\vu=-K\vx$.\\
5 & Riccati & $P_0$ & $P_\infty$ & Look one more step ahead each time until nothing changes.\\
6 & DARE & $A,B,Q,R$ & $P$, $K$ & Jump straight to that limit; $\vx_0^\top P\vx_0$ is the total score.\\
7 & MPC & model, limits, $P_f$ & a problem & Plan the next second without breaking any rule.\\
8 & QP & problem, $\vx_0$ & $H,\bm f,G,\bm h$, $\mathbf U^*$ & A bowl inside a fence; one bottom, found in milliseconds.\\
9 & Receding horizon & $\vx_k$ & $\vu_0^*$ & Use the first move only, then plan again.\\
10 & Apply to drone & $\vu_0^*$ or $-K\vx_k$ & $\Omega_1..\Omega_4$ & Add hover thrust back, share among the rotors, fly.\\
\bottomrule
\end{tabularx}
\end{center}
